%% ----------------------------------------------------------------
%% Section 2.7: Parameter-Efficient Fine-Tuning
%% ----------------------------------------------------------------
\section{Parameter-Efficient Fine-Tuning}
\label{sec:peft}

The preceding sections have established that reinforcement learning from
human feedback---and its modern variants such as GRPO---can substantially
improve LLM alignment and reasoning capabilities.  However, these
algorithms presuppose the ability to \emph{fine-tune} the underlying
language model, an operation whose cost scales linearly with the number
of trainable parameters.  For a model with $N$ parameters, full
fine-tuning requires storing not only the parameters themselves but also
their gradients and optimizer states.  With the Adam
optimizer~\citep{kingma2015adam}, this amounts to roughly $16N$ bytes of
GPU memory in mixed-precision training (2 bytes for the \texttt{fp16}
parameters, 4 bytes for the \texttt{fp32} master copy, 4 bytes each for
the first- and second-moment estimates, and 2 bytes for the gradients).
For a 7-billion-parameter model, this alone exceeds 100\,GB---well
beyond the capacity of a single commodity GPU.

The problem is compounded in reinforcement learning settings.  GRPO, for
instance, requires maintaining a \emph{reference policy} $\pi_{\text{ref}}$
alongside the \emph{active policy} $\pi_\theta$ in order to compute the
KL-divergence penalty that prevents reward hacking.  Naively, this doubles
the memory footprint attributable to model parameters.  Furthermore,
GRPO's group-based advantage estimation benefits from large group sizes
$G$ (typically 8--64 completions per prompt), which in turn demands
substantial memory for storing the sampled sequences and their
log-probabilities.  Under these constraints, full fine-tuning of models
at the 7B scale and beyond becomes impractical without access to large
multi-GPU clusters.

\emph{Parameter-efficient fine-tuning} (PEFT) methods address this
bottleneck by training only a small subset of parameters---or injecting
a small number of new trainable parameters---while keeping the bulk of
the pre-trained model frozen.  This section reviews the most prominent
PEFT techniques, with particular emphasis on Low-Rank Adaptation (LoRA)
and its quantised variant QLoRA, which together form the backbone of our
experimental pipeline.

%% ----------------------------------------------------------------
\subsection{Low-Rank Adaptation (LoRA)}
\label{sec:lora}

Low-Rank Adaptation, introduced by \citet{hu2021lora}, is grounded in the
empirical observation that the weight updates learned during fine-tuning
occupy a low-dimensional subspace relative to the full parameter space.
That is, although the pre-trained weight matrix $\mathbf{W}_0 \in
\mathbb{R}^{d \times k}$ is full-rank, the task-specific update
$\Delta\mathbf{W}$ can be well-approximated by a matrix of much lower
rank.  LoRA exploits this structure by \emph{freezing} $\mathbf{W}_0$
and injecting a trainable low-rank decomposition:
\begin{equation}
  \label{eq:lora-weight}
  \mathbf{W} = \mathbf{W}_0 + \Delta\mathbf{W}
             = \mathbf{W}_0 + \mathbf{B}\mathbf{A},
\end{equation}
where $\mathbf{B} \in \mathbb{R}^{d \times r}$ and
$\mathbf{A} \in \mathbb{R}^{r \times k}$, with the rank
$r \ll \min(d, k)$.

\paragraph{Forward pass.}
During the forward pass, the output of the adapted layer is computed as:
\begin{equation}
  \label{eq:lora-forward}
  \mathbf{h} = \mathbf{W}_0 \mathbf{x} + \mathbf{B}\mathbf{A}\mathbf{x}
             = \mathbf{W}_0 \mathbf{x}
               + \frac{\alpha}{r}\,\mathbf{B}\mathbf{A}\mathbf{x},
\end{equation}
where $\alpha$ is a scaling hyperparameter (often set equal to $r$ so that
the scaling factor $\alpha / r = 1$).  The pre-trained computation
$\mathbf{W}_0\mathbf{x}$ proceeds exactly as in the frozen model; only
the low-rank branch $\mathbf{B}\mathbf{A}\mathbf{x}$ introduces trainable
parameters.

\paragraph{Initialization.}
A key design choice in LoRA is the initialization strategy.
\citet{hu2021lora} initialize $\mathbf{A}$ with a random Gaussian
distribution, $\mathbf{A} \sim \mathcal{N}(0, \sigma^2)$, and set
$\mathbf{B} = \mathbf{0}$.  This ensures that the product
$\mathbf{B}\mathbf{A} = \mathbf{0}$ at the start of training, so the
adapted model begins from exactly the same function as the pre-trained
model.  Training therefore commences from a known-good initialization,
which stabilizes the optimization trajectory and is particularly
important when the downstream task involves reinforcement learning
objectives that are already difficult to optimize.

\paragraph{Inference-time merging.}
Unlike adapter-based methods that add sequential modules and thereby
increase inference latency, LoRA introduces \emph{no additional latency
at inference time}.  Once training is complete, the low-rank matrices can
be absorbed into the pre-trained weight:
\begin{equation}
  \label{eq:lora-merge}
  \mathbf{W}_{\text{merged}} = \mathbf{W}_0 + \mathbf{B}\mathbf{A}.
\end{equation}
The merged weight $\mathbf{W}_{\text{merged}}$ has the same dimensions as
$\mathbf{W}_0$, so the deployed model incurs identical computational cost
to the original pre-trained model.  This property is especially valuable
in production settings where serving multiple task-specific variants of
the same base model is required: one need only store and swap the small
$\mathbf{B}$ and $\mathbf{A}$ matrices rather than full model copies.

\paragraph{Rank selection.}
The rank $r$ controls the trade-off between expressiveness and efficiency.
Common choices in the literature include $r \in \{8, 16, 32, 64\}$.
\citet{hu2021lora} demonstrated that even $r = 4$ or $r = 8$ can match
full fine-tuning performance on many NLU benchmarks, suggesting that
task-specific adaptations are indeed low-rank.  For more complex tasks
such as mathematical reasoning, practitioners often increase $r$ to 16
or 32 to provide the adaptation with additional capacity.

\paragraph{Target modules.}
LoRA can, in principle, be applied to any linear layer in the
transformer.  In practice, \citet{hu2021lora} found that adapting the
\emph{query} ($\mathbf{W}_Q$) and \emph{value} ($\mathbf{W}_V$)
projection matrices in the multi-head self-attention mechanism yields
the best trade-off between parameter count and downstream performance.
Some subsequent works additionally adapt the key projection
$\mathbf{W}_K$, the output projection $\mathbf{W}_O$, and the
feed-forward network (FFN) layers, though with diminishing returns
relative to the additional parameter budget.

\paragraph{Parameter savings.}
The memory savings afforded by LoRA are dramatic.  For a single adapted
weight matrix $\mathbf{W}_0 \in \mathbb{R}^{d \times k}$, full
fine-tuning requires $d \times k$ trainable parameters, whereas LoRA
requires only $r \times (d + k)$.  For a typical transformer hidden
dimension $d = 4096$ and $r = 16$, this represents a reduction factor
of $d \times k \,/\, [r(d+k)] \approx d / (2r) = 128\times$ per
adapted matrix.  Across all attention projections in a 7-billion-parameter
model, LoRA with $r = 16$ typically trains roughly 0.1\% of the total
parameters---approximately 10 million trainable parameters out of 7
billion---while achieving performance competitive with full fine-tuning.
The corresponding reduction in optimizer state memory is proportional,
bringing the training memory footprint within reach of a single
high-end GPU.

%% ----------------------------------------------------------------
\subsection{QLoRA: Quantised Low-Rank Adaptation}
\label{sec:qlora}

While LoRA dramatically reduces the number of \emph{trainable} parameters,
the \emph{frozen} pre-trained weights $\mathbf{W}_0$ must still reside in
GPU memory for the forward and backward passes.  For a 65-billion-parameter
model stored in \texttt{float16}, this alone requires approximately 130\,GB
of memory.  \citet{dettmers2023qlora} introduced QLoRA to address this
remaining bottleneck through three complementary innovations.

\paragraph{4-bit NormalFloat quantization.}
QLoRA quantizes the frozen base model weights to 4-bit precision using
a \emph{NormalFloat} (NF4) data type.  The NF4 format is
information-theoretically optimal for normally distributed weights: its
quantization levels are chosen such that each bin contains an equal
number of values under a standard normal distribution.  Formally, for a
weight tensor $\mathbf{W}$ with entries approximately distributed as
$w_{ij} \sim \mathcal{N}(0, \sigma^2)$, the NF4 quantization levels
$q_i$ are determined by:
\begin{equation}
  q_i = \Phi^{-1}\!\left(\frac{2i + 1}{2^{b+1}}\right), \quad
  i = 0, 1, \ldots, 2^b - 1,
\end{equation}
where $\Phi^{-1}$ is the inverse of the standard normal CDF and $b = 4$
is the bit width.  During the forward pass, quantized weights are
dequantized to \texttt{bfloat16} on the fly for matrix multiplication;
the LoRA adapters $\mathbf{B}$ and $\mathbf{A}$ remain in
\texttt{bfloat16} throughout.

\paragraph{Double quantization.}
Blockwise quantization requires storing a scaling constant per block
(typically one \texttt{float32} constant per 64 weights).  These
constants themselves consume non-trivial memory at scale.  Double
quantization applies a second round of quantization to these scaling
constants, reducing their memory footprint from 32 bits to approximately
8 bits per constant.  \citet{dettmers2023qlora} report that double
quantization saves an additional 0.37 bits per parameter on average,
translating to approximately 3\,GB of savings for a 65B model.

\paragraph{Paged optimizers.}
To handle the occasional memory spikes that arise during gradient
checkpointing and dynamic batch construction, QLoRA employs
\emph{paged optimizers} that leverage NVIDIA's unified memory to
automatically page optimizer states between GPU and CPU memory.  This
mechanism prevents out-of-memory errors during training without
requiring the practitioner to manually manage memory allocation, and
introduces negligible overhead because optimizer-state access patterns
exhibit high temporal locality.

\paragraph{Practical impact.}
Taken together, these three techniques enable fine-tuning a
65-billion-parameter model on a \emph{single} 48\,GB GPU (such as an
NVIDIA A6000 or A100-48GB).  \citet{dettmers2023qlora} demonstrated that
QLoRA matches the performance of full 16-bit fine-tuning on a range of
benchmarks, establishing that the quantization-induced approximation
error is negligible for downstream task quality.  For our scaling
experiments, QLoRA is the primary method by which we fit 7B-class models
and their associated RL infrastructure (reference policy, group sampling
buffers) onto the available hardware.

%% ----------------------------------------------------------------
\subsection{Other Parameter-Efficient Methods}
\label{sec:other-peft}

While LoRA and QLoRA are the methods most directly relevant to our work,
several alternative PEFT approaches deserve brief mention, as they
provide important context for the design space.

\paragraph{Prefix tuning.}
\citet{li2021prefix} proposed prepending a sequence of trainable
continuous vectors---termed \emph{prefix tokens}---to the key and value
representations at every transformer layer.  These virtual tokens are
optimized while all pre-trained parameters remain frozen.  Prefix tuning
is effective for natural-language generation tasks but introduces
additional computation proportional to the prefix length at every layer,
slightly increasing inference latency.

\paragraph{Prompt tuning.}
\citet{lester2021prompt} introduced a simpler variant in which trainable
\emph{soft prompts}---continuous embeddings concatenated to the input
embedding sequence---are optimized for each downstream task.  Unlike
prefix tuning, soft prompts are applied only at the input layer, making
the method extremely lightweight (often fewer than 100k trainable
parameters).  However, prompt tuning tends to underperform on small
models and is less suited to complex generation tasks such as
multi-step mathematical reasoning.

\paragraph{Adapter layers.}
\citet{houlsby2019adapters} proposed inserting small bottleneck modules
(\emph{adapters}) between existing transformer layers.  Each adapter
consists of a down-projection, a nonlinearity, and an up-projection,
with a residual connection.  While adapters are highly effective and were
among the first PEFT methods proposed, they introduce sequential
computation that increases inference latency---a disadvantage that LoRA's
merging property directly addresses.

\medskip
\noindent
Table~\ref{tab:peft-comparison} summarizes the key characteristics of
these methods.

\begin{table}[t]
  \centering
  \caption{Comparison of parameter-efficient fine-tuning methods.
           ``Extra inference cost'' indicates whether the method adds
           computation beyond the original pre-trained model at
           deployment time.}
  \label{tab:peft-comparison}
  \small
  \begin{tabular}{@{}p{0.30\textwidth}p{0.20\textwidth}p{0.22\textwidth}p{0.16\textwidth}@{}}
    \toprule
    \textbf{Method} & \textbf{Trainable params} & \textbf{Extra inference cost}
      & \textbf{Merges into base} \\
    \midrule
    Full fine-tuning              & 100\%          & None       & N/A \\
    Adapters \citep{houlsby2019adapters}
                                  & $\sim$3--5\%   & Yes        & No \\
    Prefix tuning \citep{li2021prefix}
                                  & $\sim$0.1--1\% & Yes        & No \\
    Prompt tuning \citep{lester2021prompt}
                                  & $<$0.01\%      & Minimal    & No \\
    LoRA \citep{hu2021lora}       & $\sim$0.1\%    & None       & Yes \\
    QLoRA \citep{dettmers2023qlora}
                                  & $\sim$0.1\%    & None\textsuperscript{*}
                                                                & Yes \\
    \bottomrule
  \end{tabular}
  \\[4pt]
  \raggedright
  \footnotesize\textsuperscript{*}After merging and re-quantization to
  the desired serving precision.
\end{table}

%% ----------------------------------------------------------------
\subsection{LoRA for Reinforcement Learning Fine-Tuning}
\label{sec:lora-rl}

The combination of LoRA with reinforcement learning objectives such as
GRPO is not merely a matter of computational convenience; it confers
several structural advantages that make it particularly well-suited to
the RL fine-tuning setting.

\paragraph{Memory budget reallocation.}
As discussed in Section~\ref{sec:grpo} (assuming the GRPO section is
labeled accordingly), GRPO computes advantages by sampling a group of $G$
completions per prompt and evaluating each against a reward function.
The quality of the advantage estimate improves with $G$, but larger
groups require proportionally more memory for storing sequences, their
log-probabilities, and intermediate computations.  By reducing the memory
consumed by trainable parameters and their optimizer states from
	$\sim 16N$ bytes (full fine-tuning with Adam) to $\sim 16r(d+k)L$
bytes (where $L$ is the number of adapted layers), LoRA frees a
substantial portion of the GPU memory budget.  This freed memory can be
directly reallocated to support larger group sizes $G$ and larger batch
sizes, both of which improve training stability and sample efficiency.

\paragraph{Stable reference policy.}
GRPO's KL-divergence penalty
$D_{\mathrm{KL}}(\pi_\theta \| \pi_{\text{ref}})$ requires access to a
reference policy.  When using LoRA, the reference policy is simply the
frozen base model $\mathbf{W}_0$---no separate copy of the model
weights is needed.  At any point during training, the reference policy's
output can be obtained by performing a forward pass with the LoRA
adapters disabled (i.e., setting $\mathbf{B}\mathbf{A} = \mathbf{0}$).
This eliminates the need to maintain two full copies of the model in
memory, effectively halving the memory overhead associated with the
reference policy and making single-GPU GRPO training feasible for
models at the 7B scale.

\paragraph{Regularization effect.}
The low-rank constraint imposed by LoRA acts as an implicit regularizer
on the policy update.  Because the adaptation $\Delta\mathbf{W} =
\mathbf{B}\mathbf{A}$ is restricted to a rank-$r$ subspace, the policy
cannot deviate arbitrarily from $\pi_{\text{ref}}$, which complements
the explicit KL penalty in GRPO.  Empirically, this dual regularization
(low-rank constraint plus KL penalty) has been observed to reduce reward
hacking---the phenomenon where the policy exploits reward model
weaknesses rather than genuinely improving on the target
task~\citep{hu2021lora, dettmers2023qlora}.

\paragraph{Efficient multi-task and iterative training.}
In iterative GRPO pipelines where the policy is periodically updated and
the reward function or curriculum changes, LoRA enables efficient
checkpointing and rollback.  Each checkpoint consists only of the small
$\mathbf{B}$ and $\mathbf{A}$ matrices (typically tens of megabytes
rather than tens of gigabytes), facilitating rapid experimentation with
different hyperparameter schedules, reward shaping strategies, and
curriculum orderings---all of which are critical dimensions in the
scaling study presented in later chapters.

\medskip
\noindent
In summary, parameter-efficient fine-tuning---and LoRA/QLoRA in
particular---is not merely an engineering convenience but a
\emph{methodological enabler} for the class of RL-based training
experiments we pursue.  By compressing the trainable parameter space,
eliminating the need for a separate reference model copy, and freeing
GPU memory for the large sample groups that GRPO requires, LoRA
transforms the scaling study from one that would demand a multi-node
GPU cluster into one achievable on modest hardware.  The specific
LoRA configurations and QLoRA quantization settings used in our
experiments are detailed in Chapter~4.
